\section{A fixed-initial-data realization by integer Fourier shears}\label{sec:parabolic}

The lower bounds at a prescribed time do not determine what happens over
a time interval.  This section treats one explicit family with its
coefficients specified at time zero.  The heat operator has a positive
time parameter throughout.

Let $T_m$ be the Chebyshev polynomial of the first kind, defined by
$T_m(\cos\theta)=\cos(m\theta)$.  For $A>0$ set
\begin{equation}\label{eq:heat-edge-definition}
 Q_{m,A}(s,c)=
 \left(e^{sA^{-2}D^2}T_m\right)(1+c/A),
 \qquad \mu=(m^2/A)^{1/3}.
\end{equation}
The exponential is a finite sum when applied to a polynomial.
All limits in this section are joint limits satisfying
\begin{equation}\label{eq:edge-degree-regime}
 A\longrightarrow\infty,\qquad
 m/\sqrt A\longrightarrow\infty,\qquad m/A^2\longrightarrow0.
\end{equation}
Equivalently, $\mu\to\infty$ and $\mu/A\to0$.
No convergence rate in these two conditions is imposed.

\begin{theorem}[Two scales in the heat-evolved edge]\label{thm:heat-edge}
Under \eqref{eq:edge-degree-regime}, for every compact interval
$S\subset(0,\infty)$ and every $C_0<\infty$, $Q_{m,A}(s,c)$ is positive
for all sufficiently large parameters, uniformly for $s\in S$ and
$|c|\le C_0$.  Moreover,
\begin{align}
 \mu^{-2}\log Q_{m,A}(s,0)&\longrightarrow
 F_s:=3\,2^{-4/3}s^{1/3},\label{eq:heat-leading-limit}\\
 \mu^{-1}\log\frac{Q_{m,A}(s,c_2)}{Q_{m,A}(s,c_1)}
 &\longrightarrow G_s(c_2-c_1),
 \qquad G_s:=2^{-2/3}s^{-1/3}.
 \label{eq:heat-tilt-limit}
\end{align}
Both convergences are uniform for $s\in S$ and $|c_1|,|c_2|\le C_0$.
In particular $F_4=3\,2^{-2/3}$ and $G_4=2^{-4/3}$.
\end{theorem}

\begin{proof}
This is the first-kind case of Theorem~\ref{thm:quantitative-edge}.
Its sufficient parameter test holds eventually along every sequence
in \eqref{eq:edge-degree-regime}; both stated error bounds tend to zero.
In particular the displacement limit uses the separate Gaussian-tilt
estimate, not subtraction of amplitude estimates.
\end{proof}

\begin{lemma}[A majorant and a terminal time interval]\label{lem:heat-terminal}
Put $D_k=T_m^{(k)}(1)$.  Then
\begin{equation}\label{eq:heat-positive-series}
 Q_{m,A}(s,c)=\sum_{\ell+2j\le m}
 \frac{D_{\ell+2j}}{A^{\ell+2j}\ell!j!}c^\ell s^j,
 \qquad
 \frac{D_{k+1}}{D_k}=\frac{m^2-k^2}{2k+1}\quad(0\le k<m).
\end{equation}
All $D_k$ in this range are positive.  For $s$ in a fixed compact
subinterval of $(0,\infty)$ and $0<c_0\le c\le C_0$,
\begin{equation}\label{eq:heat-log-derivatives}
 0\le\partial_s\log Q\le C(1+\mu^2),\qquad
 0\le\partial_c\log Q\le C(1+\mu).
\end{equation}
Thus changes of size at most $h_\mu=(1+\mu^2)^{-1}$ in $s$ and $c$
change $Q$ by at most a fixed factor on these compact sets.
If $d\ge0$, $0\le s\le S_1$, and $|x|\le1+d/A$, then
\begin{equation}\label{eq:heat-global-majorant}
 \left|\left(e^{sA^{-2}D^2}T_m\right)(x)\right|
 \le Q_{m,A}(S_1,d).
\end{equation}
\end{lemma}

\begin{proof}
Differentiate $(1-x^2)T_m''-xT_m'+m^2T_m=0$ $k$ times and evaluate
at $x=1$.  This gives the displayed recurrence, starting from $D_0=1$.
The finite Taylor expansion followed by the finite heat expansion gives
\eqref{eq:heat-positive-series}.  At fixed $\ell$, the ratio of the
$(j+1)$st and $j$th nonzero summands is at most
\[
 \frac{s m^4}{A^2(j+1)(4j+1)(4j+3)}
 \le\frac{C\mu^6}{(j+1)^3}.
\]
For $j\ge J=\lceil C_1(1+\mu^2)\rceil$ this is at most $1/2$.
Under the normalized weights in that row, the tail above $J$ is bounded
by a geometric sequence, so its mean index is at most $J+2$.
Rows ending before $J$ satisfy the same bound, and zero rows are ignored.
Averaging over rows gives
$\partial_s\log Q=\mathbb E(j)/s\le C(1+\mu^2)$.
The identity $\partial_c^2Q=\partial_sQ$ yields
$\mathbb E[\ell(\ell-1)]=c^2\partial_sQ/Q$.
Since $\ell\le1+\sqrt{\ell(\ell-1)}$, Cauchy--Schwarz gives
$\partial_c\log Q\le c^{-1}+\sqrt{\partial_sQ/Q}$.
Integrating these derivative estimates proves the assertion about
$h_\mu$.

For completeness, $|T_m^{(k)}(x)|\le D_k$ for $|x|\le1$.
For $k\ge1$ this follows from
$T_m^{(k)}=m2^{k-1}(k-1)!C_{m-k}^{k}$: the generating function
\[
 \sum_{n\ge0}C_n^k(\cos\theta)t^n
 =(1-te^{i\theta})^{-k}(1-te^{-i\theta})^{-k}
\]
has nonnegative absolute coefficient majorants given by its value at
$\theta=0$.  The case $k=0$ is the defining Chebyshev bound.
For $1\le|x|\le1+d/A$, parity and the positive Taylor coefficients at
one instead give
$|T_m^{(k)}(x)|\le T_m^{(k)}(1+d/A)$.
The heat series has nonnegative coefficients in $s$, so termwise use of
these derivative majorants proves \eqref{eq:heat-global-majorant}.
\end{proof}

\subsection{Uniform transfer to consecutive integer modes}

The following estimate keeps the polynomial degree in the error bound.
This is necessary because pointwise convergence at fixed degree would
not justify the subsequent limit.

\begin{proposition}[Dynamic confluence]\label{prop:dynamic-confluence}
Fix $\delta_0,\beta>0$ and let
\begin{equation}\label{eq:parabolic-frozen-parameters}
 a=\frac{32}{63}L,\quad A=a-\delta_0,\quad
 R=e^{-9L^2/40000},\quad e_a=A/a,\quad
 v=-\beta/A,\quad B=av/R.
\end{equation}
Suppose $m$ is a positive even integer satisfying
\eqref{eq:edge-degree-regime}.  Set $\eta=AR$ and
\begin{equation}\label{eq:parabolic-fourier-coefficients}
 w_\eta(x)=\frac{e^{i\eta x}-1}{i\eta},\qquad
 T_m(w_\eta(x))=\sum_{j=0}^m b_j(\eta)e^{ij\eta x},\qquad
 b_j(\eta)=\frac{T_m^{(j)}(i/\eta)}{j!(i\eta)^j}.
\end{equation}
Every $b_j$ is nonzero.  Define the real periodic solution
\begin{equation}\label{eq:parabolic-exact-shear}
 F(t,y)=2\Rea\sum_{j=0}^m b_j(\eta)
 e^{-(m+j)^2t+i(m+j)(y-Bt)},\qquad
 u(t,x)=(0,B,F(t,x_2)).
\end{equation}
Then $u$ is a smooth unforced unit-viscosity Navier--Stokes solution with constant
pressure on $(\mathbb R/2\pi\mathbb Z)^3$ and exactly $m+1$ positive
frequencies $m,\ldots,2m$ in its third component.
Writing $G(s,z)=F(R^2s,aRz)$ and
$P_s(x)=(e^{sA^{-2}D^2}T_m)(x)$, one has
\begin{equation}\label{eq:parabolic-uniform-transfer}
 \sup_{0\le s\le4,\ |x|\le2}
 \left|G(s,e_ax)-2P_s(x-vs/e_a)\right|
 \le \exp\!\left(-\frac9{40000}L^2+o(L^2)\right).
\end{equation}
\end{proposition}

\begin{proof}
The coefficient identity is the Taylor formula for $T_m$ at $i/\eta$.
Every derivative of a real-rooted polynomial is real-rooted, by Rolle's
theorem, and $T_m$ has only real roots.  Thus $T_m^{(j)}(i/\eta)\ne0$
for $j<m$; the derivative of order $m$ is a nonzero constant.
The equation $(\partial_t+B\partial_y-\partial_y^2)F=0$ follows directly
from the finite sum.  It gives the claimed divergence-free shear solution.

Put $M=m\eta$ and $t_s=s/A^2$.  The complex half of $G(s,e_ax)$ is
exactly
\begin{equation}\label{eq:parabolic-conjugation}
 H(s,x)=e^{-M^2t_s+iM(x-vs/e_a)}
 \left(e^{t_sD^2}T_m(w_\eta(\cdot))\right)
 (x-vs/e_a+2iMt_s).
\end{equation}
This identity can be checked on every exponential in
\eqref{eq:parabolic-fourier-coefficients}; in particular the imaginary
shift has a positive sign.  Under the change of variable $w=w_\eta(x)$,
the exact polynomial generator is
\[
 \mathcal L_\eta=(1+i\eta w)^2D_w^2
                  +i\eta(1+i\eta w)D_w.
\]
It maps polynomials of degree at most $m$ to the same space.
The chain rule gives, for every polynomial $p$,
\[
 D_x^2(p\circ w_\eta)=(\mathcal L_\eta p)\circ w_\eta,
 \qquad
 e^{tD_x^2}(p\circ w_\eta)
       =(e^{t\mathcal L_\eta}p)\circ w_\eta .
\]
The second identity follows from the first by the exponential series:
both operators act on the same finite-dimensional invariant space,
and the matrix series converge at every finite time.
In the norm $\|p\|_4=\sum_{k=0}^m|p_k|4^k$ on that space,
\[
 \|\mathcal L_\eta-D_w^2\|\le C\eta m^2,
 \qquad \|\mathcal L_\eta\|+\|D_w^2\|\le Cm^2
\]
for $0<\eta\le1$.  These bounds follow termwise from
$\|D_w\|\le m/4$ and $\|D_w^2\|\le m^2/16$.
The recurrence $T_{m+1}=2wT_m-T_{m-1}$ gives
$\|T_m\|_4\le C_4^m$.  Duhamel's formula therefore gives
\[
 \|e^{t_s\mathcal L_\eta}T_m-e^{t_sD_w^2}T_m\|_4
 \le C t_s\eta m^2 C_4^m e^{Ct_sm^2}.
\]
For $|x|\le2$, all the complex arguments in
\eqref{eq:parabolic-conjugation} have modulus less than $5/2$ for large
$L$, and their $w_\eta$ images have modulus less than $3$.
Thus they remain inside the coefficient disk of radius four.
On these smaller disks $|w_\eta(z)-z|\le C\eta$, while one derivative
of a polynomial costs at most $m/4$ in the coefficient norm.
Replacing the generator, replacing $w_\eta(z)$ by $z$, removing the
imaginary shift $2iMt_s$, and removing the carrier and scalar factor
in \eqref{eq:parabolic-conjugation} consequently give a total error at
most
\[
 \eta\,\operatorname{poly}(m,A)
       \exp(Cm+Cm^2/A^2).
\]
The real drift $vs/e_a$ has been retained in the comparison function.
Now $\log\eta=-9L^2/40000+O(\log L)$,
$m=o(A^2)=o(L^2)$, and $m^2/A^2=o(L^2)$.
Taking twice the real part proves \eqref{eq:parabolic-uniform-transfer}.
\end{proof}

\begin{remark}
Every fixed member of \eqref{eq:parabolic-exact-shear} has finite energy
on the torus.  The coefficients and the constant drift need not remain
bounded as $L\to\infty$.  The periodic extension to $\mathbb R^3$ has
infinite total energy and is used only in local coordinates.  No uniform
energy class or limiting weak solution is asserted here.
\end{remark}

\subsection{The signed flux and the whole spatial plateau}

Fix $0<\delta_0<\delta$ and a nonnegative smooth function $\vartheta$
supported in $(-\delta,\delta)$ that equals one on
$[-\delta_0,\delta_0]$.  Suppose there are $r_c,h,\vartheta_*>0$ such that
\begin{equation}\label{eq:parabolic-outer-cap}
 \delta_0<r_c-2h<r_c+2h<\delta,
 \qquad \vartheta(r)\ge\vartheta_*
 \quad\hbox{for }|r-r_c|\le2h.
\end{equation}
Let $\chi\in C^1([0,4])$ satisfy $0\le\chi\le1$, $\chi(0)=0$, and $\chi=1$ on a
fixed neighborhood of four.  Define the annular cutoff and its plateau by
\[
 \phi_L(x)=\vartheta(|x|/R-a),\qquad
 \mathcal P_L=\{x:(a-\delta_0)R\le|x|\le(a+\delta_0)R\}.
\]
For large $L$ these sets lie inside a Euclidean coordinate neighborhood
of the torus.  For the shear of Proposition~\ref{prop:dynamic-confluence}
put
\begin{align}
 \mathcal T_L&=\frac12\int_0^{4R^2}\chi(t/R^2)
  \int_{\mathbb T^3}|u(t,x)|^2u(t,x)\cdot\nabla\phi_L(x)\,dx\,dt,
 \label{eq:parabolic-physical-flux}\\
 M_L^{\mathrm{plat}}&=\int_0^{4R^2}\int_{\mathcal P_L}
              |F(t,x_2)|^3\,dx\,dt.
 \label{eq:parabolic-physical-payment}
\end{align}
The denominator integrates the oscillatory velocity component over the
entire spatial plateau and the entire time interval.  It is not the
$L^3$ mass of the full velocity: the constant drift $B$ is not included
in \eqref{eq:parabolic-physical-payment}.

\begin{theorem}[A full-interval residual for one exact shear family]
\label{thm:parabolic-shear}
Use the parameters and shear of Proposition~\ref{prop:dynamic-confluence}
and the fixed cutoffs above.  Write
\[
 c_-=r_c-h+\delta_0,\qquad c_p=2\delta_0,\qquad
 \Delta_c=c_--c_p=r_c-h-\delta_0>0.
\]
There are constants $c,C>0$ depending only on the fixed cutoffs and
$\beta$ such that, along every degree sequence
\eqref{eq:edge-degree-regime}, eventually $\mathcal T_L>0$ and
\begin{align}
 \frac{c\beta a^{2/3}}{A(1+\mu^2)}
  \exp\!\left((2^{-1/3}\Delta_c+o(1))\mu\right)
 &\le R^{1/3}\frac{\mathcal T_L}{(M_L^{\mathrm{plat}})^{2/3}}
 \label{eq:parabolic-residual-lower}\\
 &\le\frac{C\beta a^{4/3}}A(1+\mu^2)^{2/3}e^{C\mu}.
 \label{eq:parabolic-residual-upper}
\end{align}
The eventual threshold may depend on the chosen degree sequence.
If in addition $\mu/\log A\to\infty$, the middle quotient diverges.
With
\[
 \omega=e^{-2L^2/3969},\qquad
 p_L=R^{-2}\omega M_L^{\mathrm{plat}},\qquad
 X_L=(\omega/R)\max\{\mathcal T_L,0\},
\]
one nevertheless has
\begin{equation}\label{eq:parabolic-normalized-rate}
 \lim_{L\to\infty}\frac1{L^2}\log\frac{X_L}{p_L^{2/3}}
 =-\frac2{11907}.
\end{equation}
\end{theorem}

\begin{proof}
The argument retains both opposed caps in the numerator and the full
plateau in the denominator.

\emph{Exact cross sections.}
For $G(s,z)=F(R^2s,aRz)$ define
\[
 W_a(z)=-2\pi az\vartheta(a(|z|-1)),
\]
\[
 \mathcal A_a(z)=\pi\left(
  [(a+\delta_0)^2-a^2z^2]_+
  -[(a-\delta_0)^2-a^2z^2]_+\right),
\]
where $[r]_+=\max(r,0)$, and set
\[
 \mathcal S=v\int_0^4\chi(s)\int_{\mathbb R}
                  W_a(z)|G(s,z)|^2\,dz\,ds,
 \qquad
 \mathcal K=\int_0^4\int_{\mathbb R}
                  \mathcal A_a(z)|G(s,z)|^3\,dz\,ds.
\]
Integrating $\partial_2\phi_L$ on the transverse plane gives
\[
 \int_{\mathbb R^2}\partial_2\phi_L(x_1,y,x_3)\,dx_1dx_3
   =-2\pi y\vartheta(|y|/R-a).
\]
Indeed polar integration of the radial derivative reduces this integral
to $2\pi y\int_{|y|}^\infty\partial_r
\vartheta(r/R-a)\,dr$.  The term involving $F\partial_3\phi_L$ in
\eqref{eq:parabolic-physical-flux} integrates to zero, and so does the
$B^3\partial_2\phi_L$ term.  The remaining changes of variables
$t=R^2s$, $y=aRz$, together with $B=av/R$, give the exact relations
\begin{equation}\label{eq:parabolic-dimensional-identities}
 \mathcal T_L=\frac{a^2R^3}{2}\mathcal S,
 \qquad M_L^{\mathrm{plat}}=aR^5\mathcal K.
\end{equation}
The second relation uses the full difference of disks represented by
$\mathcal A_a$, including its outer caps.

\emph{Paired cap sign.}
Let $\gamma_L=\beta a/A$, so $\gamma_L\to\beta$.  At
$z=\pm(1+r/a)$ the evenness of $P_s$ in
Proposition~\ref{prop:dynamic-confluence} gives the ideal values
\[
 2Q_{m,A}(s,r+\delta_0+\gamma_Ls),\qquad
 2Q_{m,A}(s,r+\delta_0-\gamma_Ls),
\]
respectively.  Theorem~\ref{thm:heat-edge}, uniformly for
$3\le s\le4$ and $r$ in the fixed support of $\vartheta$, implies
\[
 \frac1\mu\log
 \frac{Q_{m,A}(s,r+\delta_0+\gamma_Ls)}
      {Q_{m,A}(s,r+\delta_0-\gamma_Ls)}
   =2G_s\gamma_Ls+o(1).
\]
The right side has a fixed positive lower bound.  Both ideal values
are positive and grow at the $\exp(F_s\mu^2)$ scale, so the absolute
error in \eqref{eq:parabolic-uniform-transfer} is negligible relative
to them.  Hence the squared exact value at the negative cap is at most
$e^{-c_1\beta\mu}$ times the squared value at the positive cap, for
some $c_1>0$ and all sufficiently large parameters.
Since $W_a$ is odd, $W_a(z)\le0$ for $z\ge0$, and $v<0$, the sum of
each paired contribution is nonnegative throughout $3\le s\le4$.

\emph{The terminal contribution and earlier times.}
Let $h_\mu=(1+\mu^2)^{-1}$.  On
\[
 4-h_\mu\le s\le4,\qquad r_c-h\le r\le r_c+h,
\]
the time cutoff is one for sufficiently large parameters and
$-W_a(1+r/a)\ge c_2 a$.  The positive ideal coordinate is at least
$c_-+4\gamma_L-C h_\mu$.  The derivative estimates in
Lemma~\ref{lem:heat-terminal} give
\[
 |G(s,1+r/a)|\ge c_3 Q_{m,A}(4,c_-+4\gamma_L).
\]
The loss from its paired cap tends to zero.  Using $dz=dr/a$ therefore
gives a contribution at least
$c_4(\beta/A)h_\mu Q_{m,A}(4,c_-+4\gamma_L)^2$.

All real arguments arising from the fixed collar on $0\le s\le3$
have absolute value at most $1+C_1/A$ for a fixed $C_1>0$.
The global majorant \eqref{eq:heat-global-majorant} and the transfer
bound thus bound the absolute contribution of this whole earlier
interval by
$C_2(\beta/A)Q_{m,A}(3,C_1)^2$.
Theorem~\ref{thm:heat-edge} gives
\[
 \log\frac{Q_{m,A}(4,c_-+4\gamma_L)}{Q_{m,A}(3,C_1)}
   =(F_4-F_3)\mu^2+o(\mu^2),\qquad F_4>F_3.
\]
Thus the ratio of the early absolute bound to the terminal lower bound
is at most
\[
 C h_\mu^{-1}\exp(-2(F_4-F_3)\mu^2+o(\mu^2)),
\]
which tends to zero.
The same global majorant on the full interval gives an upper bound.
Consequently, for one fixed $C_+>0$,
\begin{equation}\label{eq:parabolic-S-bounds}
 0<c\frac\beta A h_\mu Q_{m,A}(4,c_-+4\gamma_L)^2
 \le\mathcal S
 \le C\frac\beta A Q_{m,A}(4,C_+)^2.
\end{equation}
The constants here use
$\int|W_a(z)|\,dz\le C$, as follows by setting
$z=\pm(1+r/a)$ on its fixed-width radial support.

\emph{Upper and lower payment on the full plateau.}
In the coordinate $z=e_ax$ the plateau projects to
$|x|\le1+2\delta_0/A$.  Transport changes the ideal real coordinate
by at most $4\gamma_L/A$.  Therefore
\eqref{eq:heat-global-majorant} and the transfer estimate give
\[
 \sup_{0\le s\le4,\ \mathcal A_a(z)>0}|G(s,z)|
 \le C Q_{m,A}(4,c_p+4\gamma_L).
\]
Since $\mathcal A_a\le4\pi a\delta_0$ and its support has bounded
length in $z$,
\begin{equation}\label{eq:parabolic-K-upper}
 \mathcal K\le Ca Q_{m,A}(4,c_p+4\gamma_L)^3.
\end{equation}
For a lower bound, restrict only the integration to the outer strip
\[
 z=1+(c-\delta_0)/a,
 \qquad \delta_0\le c\le\tfrac32\delta_0,
 \qquad 4-h_\mu\le s\le4.
\]
Here $\mathcal A_a(z)\ge c_5a$ and $dz=dc/a$.
The positive real field, the transfer estimate, and
Lemma~\ref{lem:heat-terminal} yield
\begin{equation}\label{eq:parabolic-K-lower}
 \mathcal K\ge c h_\mu
       Q_{m,A}(4,\delta_0+4\gamma_L)^3.
\end{equation}
Only this lower bound restricts to a strip; the denominator itself in
\eqref{eq:parabolic-physical-payment} remains the full plateau integral.

By \eqref{eq:heat-tilt-limit}, uniformly for the varying but bounded
arguments here,
\[
 \log\frac{Q_{m,A}(4,c_-+4\gamma_L)}
              {Q_{m,A}(4,c_p+4\gamma_L)}
     =(2^{-4/3}\Delta_c+o(1))\mu,
 \qquad
 \log\frac{Q_{m,A}(4,C_+)}{Q_{m,A}(4,\delta_0+4\gamma_L)}
     \le C\mu.
\]
Now \eqref{eq:parabolic-dimensional-identities} gives
\[
 R^{1/3}\frac{\mathcal T_L}{(M_L^{\mathrm{plat}})^{2/3}}
     =\frac{a^{4/3}}2\frac{\mathcal S}{\mathcal K^{2/3}}.
\]
Combining the lower bound for $\mathcal S$ with
\eqref{eq:parabolic-K-upper} proves
\eqref{eq:parabolic-residual-lower}; combining the upper bound with
\eqref{eq:parabolic-K-lower} proves
\eqref{eq:parabolic-residual-upper}.  Their polynomial factors have
logarithms $O(\log A)$ because $\mu=o(A)$, so the stated divergence
under $\mu/\log A\to\infty$ follows.

Finally the exact normalization is
\[
 \frac{X_L}{p_L^{2/3}}
  =\omega^{1/3}R^{1/3}
       \frac{\mathcal T_L}{(M_L^{\mathrm{plat}})^{2/3}}.
\]
Both logarithmic bounds for the second factor are
$O(\mu+\log A)=o(L^2)$, whereas
$L^{-2}\log\omega^{1/3}=-2/11907$.
This proves \eqref{eq:parabolic-normalized-rate}.
\end{proof}

\begin{remark}[Scope of the exponent]
At a fixed slice the Chebyshev edge exponent is of size
$m/\sqrt A=\mu^{3/2}$.  Theorem~\ref{thm:parabolic-shear} gives a
residual exponential factor of size $\exp(\Theta(\mu))$ after time
integration and division by the full plateau payment, with the explicit
polynomial factors displayed in the theorem.  Describing the logarithm
of the entire quotient as $\Theta(\mu)$ additionally requires
$\log A=o(\mu)$.  The negative normalized rate applies to this one
family, with its specified coefficients, drift, and degree range.
It is not a uniform estimate over all packets and gives no regularity or
singularity conclusion for general Navier--Stokes solutions.
\end{remark}

\begin{remark}[Degrees outside the proof]
When $m$ is comparable to or larger than $A^2$, the relation
$\mu/A\to0$ fails and the edge saddle used above is no longer the
relevant expansion.  In addition, the error estimate in
Proposition~\ref{prop:dynamic-confluence} contains
$\exp(Cm+Cm^2/A^2)$.  It does not tend to zero from the argument given
here in such a range.  In particular a formal calculation at
$m\asymp A^4$ is not a consequence of any theorem in this paper.
\end{remark}
